\clearpage
\section{Audit and verification record for v0.10.0}\label{app:audit-v0100}

\subsection{Source authority and evidence status}
The source base is the supplied v0.9.0 manuscript and full package. Their SHA-256 hashes are recorded in \texttt{audit/base\_sha256.json}, and the three original files are preserved in \texttt{previous/}. This continuation adds K62--K71. All inherited theorem, lemma, proposition, corollary, and proof blocks are preserved verbatim. Current prose and the ledger now distinguish the refuted exchange-converse claim from the proved fixed-finite-group lifting criterion. Historical version and audit records remain historical.

The new block has written proofs and an authoring-session self-review, not independent review or Lean certification. It does not assert that all earlier mathematics received a fresh audit. The Lean and Lake executables were not available on the runtime path. This is a tool limitation, not a mathematical obstruction. No historical novelty is claimed for the deformation-theoretic or finite-group stability results.

\subsection{Executed routes and decisive outcomes}
Three routes were pursued. First, the normalized multiplicativity defect gives a complete next-level obstruction in $H^2(G,\operatorname{Ad}\bar\rho)$, with explicit linear repair and nonexistence certificates. Second, a central sign factor makes the derived-invariant test blind to an action that still cannot lift. This refutes sufficiency of that one diagnostic. Third, finite averaging gives a convergent correction whenever defect depth exceeds twice the valuation of the group order. Its finite-horizon corollary supplies a necessary-and-sufficient existence test for a fixed finite group.

The corresponding ultrametric stability threshold is established in \cite[Lemma 7.3]{FournierFacio2026}. The source was checked at the exact v5 version, with the displayed hypotheses and proof, not from its abstract. The direct matrix proof here is independent of that citation and supplies the implementation contract. The normalized group-cohomology conventions were checked against \cite[Chapter II, 1.17--1.19]{MilneCFT2020}. The new source record states the translation and the limits of the overlap check. An initial broad search returned irrelevant material, which was discarded. No inference of novelty is drawn from that search.

\subsection{Proof boundaries examined}
\begin{enumerate}
\item The ideal $p^nR_{n+1}$ is square-zero even at $n=1$. Matrix corrections use left multiplication, which fixes the signs in the cocycle equation and strict-conjugacy action.
\item The next-level obstruction is a class, not the chosen raw cocycle. The affine equation is necessary and sufficient for the entire fixed group multiplication table. A failed choice of representative matrices is not by itself an obstruction.
\item The approximate conjugations $\alpha_g$ are not an action. The averaging proof starts with an exact associativity identity and controls its quadratic remainder before averaging. Division by $|G|$ loses $e=v_p(|G|)$ digits, explicitly accounted for at each step.
\item The convergence threshold is strictly $m>2e$. The finite witness must extend the requested $\rho_n$, but the reconstructed integral action need only preserve that level, not the witness's top $e$ digits. The example $3\pmod8$ for $C_2$ verifies that this distinction is real.
\item The central-sign example uses odd $p$, so its idempotents are defined. It is a counterexample for a fixed finite group; its restriction to $C_p$ supplies the impossible relation. No identification of its entire exchange cone with the adjoint $H^2$ class is made.
\item Finite catalogue means a finite set of isomorphism classes. It does not claim finite automorphism groups or a category equivalence under reduction. Strict framed matrices need not be unique.
\item The profinite example changes the source from $C_p$ to $\mathbb Z_p$. The latter has an integral lift because it does not impose the same finite-order relation. This prevents applying a fixed-finite-quotient nonlifting verdict indiscriminately to a Galois representation.
\end{enumerate}

\subsection{Finite computation contract}
All five exact regression scripts passed, including reruns of the four inherited suites. The new script uses integer and modular arithmetic only in mathematical checks. It examines eight specified cochain fixtures, including the nonabelian group $S_3$, with 593 cocycle triples and 56 changes of section. Positive solutions and negative left-nullspace witnesses are checked directly against their equations.

The cyclic suite enumerates all 367 valid rank-two actions and all 40,987 next-level generator matrices in its six explicitly listed $(p,n,q)$ cases. The averaging suite performs 330 corrections in 110 deterministic fixtures, including primes dividing and not dividing the group order. Every use of division tracks finite working precision, and the exact associativity identity is checked before the predicted new defect depth. There are also 36 rank-one horizon comparisons, 98 conjugacy-average checks, and all 3,107 next-level lifts of three selected non-scalar nonliftable matrices. The sign and profinite power calculations are supplementary finite checks.

These are bounded tests of the implementation. They do not establish $p$-adic convergence, exhaustive classification for arbitrary groups, or historical novelty. Those conclusions either have their written proofs above or are not claimed. Full iterators, exclusions, counts, interpreter version, wall times, command lines, hashes before and after execution, and report hashes are stored in the package. Each suite has a 40-second subprocess timeout. No memory cap was set. Numerical-library thread limits are cooperative.

\subsection{Manuscript and continuation}
The standalone source is reproduced by \texttt{audit/integrate\_v0100.py}. The source diff, inheritance checks, LaTeX log, rendering record, and checksum manifest are included. No external source text or font files are bundled. Final proofreading checks the new formulas, hypotheses, labels, references, and surrounding integration prose.

The sufficient existence question is closed here only for a fixed finite source group over $\mathbb Z_p$. Continuous lifting with growing finite images, geometric local conditions, and automorphic realization remain separate obligations. A useful continuation must address those actual hypotheses rather than merely rename the already-classical finite-group obstruction or averaging argument.
